\section{Conclusion}
\label{sec:conclusion}
This paper presents a reproducible diagnostic protocol for sparse-concept and calibration evaluation in Knowledge Tracing. The protocol combines learner-based and temporal splits, train-only KC-frequency stratification with an explicit strict cold-start group, per-stratum sample-size-aware reporting (ECE, Brier decomposition, reliability diagrams), and an eight-channel leakage and predictive-sanity audit. 

Our experimental application of the protocol reveals several key findings. First, while aggregate population metrics may remain competitive, model calibration tends to degrade as KC training evidence decreases; for instance, SimpleKT's Expected Calibration Error (ECE) on ASSISTments 2012 reliably degraded from $0.113$ on dense concepts to $0.158$ on medium concepts, and further to $0.225$ on sparse concepts (interpreted with its Limited reliability flag). Second, performance and calibration trends exhibit pronounced heterogeneity across datasets, underscoring the necessity of dataset-local diagnostics. Finally, our temporal-split audit demonstrated that the L8 channel successfully catches subtle evaluation pipeline artifacts, such as prediction--label misalignment, ensuring that researchers can distinguish genuine concept-level cold-start failures from implementation errors.

Future work can extend this diagnostic foundation by evaluating across multiple temporal seeds or cutoffs to capture variance, by incorporating a classical baseline that exploits item difficulty for unseen learners under learner-based splits, and by applying the protocol to diagnose graph-augmented and self-supervised KT architectures.
